\section{Introduction}\label{sec:introduction}

An exponentially long sum can have a short arithmetic computation.
Choose one variable from each of the \(n\) pairs \((x_i,y_i)\) and
multiply the chosen variables. There are \(2^n\) choices. Their
products sum to \(\prod_{i=1}^n(x_i+y_i)\).
Now suppose a polynomial-size arithmetic circuit takes a string of
polynomially many bits as an additional input and computes the term
indexed by that string. Can the sum over all strings always be
computed by a polynomial-size circuit? Valiant's class
\(\mathsf{VNP}\) allows this summation~\cite{Valiant}. Membership in
\(\mathsf{VP}\) requires a polynomial-size circuit for the resulting
polynomial itself. We prove that \(\mathsf{VP}\ne\mathsf{VNP}\) over
every field.

Our witness is a polynomial on a square grid. Give each edge a label
and place a table of independent variables at each vertex. A labeling
of the grid selects one entry at each vertex according to the labels
on the edges meeting there. The product of those entries is the
weight of the labeling. On an \(n\times n\) grid with \(n\) possible
labels per edge, a polynomial-size circuit can compute this weight
from the labels. Summing the weights over all labelings gives the
grid's \emph{state sum}. We show that its circuit size grows
superpolynomially with \(n\), even allowing cancellations and
unrestricted reuse of intermediate values.

The same construction on two simpler graphs separates the models
within \(\mathsf{VP}\). On a complete binary tree, a polynomial-size
circuit can compute and reuse the contributions from its subtrees.
Algebraic branching programs (ABPs) require superpolynomial size for
this family. On a path, the state sum is an entry of a product of
matrices. A branching program can evaluate it by remembering the
state of the current edge, and here the superpolynomial lower bound
is for arithmetic formulas. Over a field \(F\), we write \(\VF\)
and \(\VBP\) for the polynomial-size formula and branching-program
classes. In all four classes the number of variables and the output
degree are polynomially bounded in the family index.

\needspace{15\baselineskip}
\begin{theorem}[Three separations]\label{thm:main}
Over every field \(F\),
\[
             \VF\subsetneq\VBP\subsetneq\VP\subsetneq\VNP .
\]
With \(n\) states per edge, the following graph families witness the
three strict inclusions. The lower bounds hold in degree-component
border complexity, with absolute constants for all sufficiently large \(n\).
\begin{center}
\small\normalfont
\begin{tabular}{@{}ll@{\hspace{1.15em}}l@{}}
\toprule
Graph & Upper class & Lower bound \\
\midrule
\(n\times n\) square grid & \(\VNP\) & Circuit size \(n^{\Omega((\log n)^{1/12})}\)\\
Complete binary tree, \(\lfloor\log_2(2n)\rfloor\) levels &
\(\VP\) & ABP size \(n^{\Omega(\sqrt{\log\log n})}\)\\
Path on \(n\) vertices & \(\VBP\) & Formula size \(n^{\Omega(\log\log n)}\)\\
\bottomrule
\end{tabular}
\end{center}
The determinant of an \(n\times n\) matrix of independent variables also
has degree-component border formula size at least \(n^{\Omega(\log\log n)}\).
\end{theorem}

To write the state sum explicitly, let \(G=(V,E)\) be a finite loopless multigraph
and let \(E(v)\) be the list of edges at \(v\). For each assignment
\(\alpha\in[q]^{E(v)}\), where \([q]=\{1,\ldots,q\}\), let
\(X_{v,\alpha}\) be the corresponding table entry. With \(q\) states
per edge, the polynomial is
\begin{equation}\label{eq:graph}
       Z_{G,q}(X)=\sum_{a\in[q]^E}\prod_{v\in V}X_{v,a|_{E(v)}} .
\end{equation}
For the grid, the \(\VNP\) description works in every characteristic.
A binary string encodes the labels. Boolean indicators pick out the
corresponding table entry at each vertex, and unused encodings
contribute zero. Since at most four edges meet at a vertex, there are
at most \(n^4\) entries to choose from in each table. The circuit lower
bound allows arbitrary constants, affine inputs, and limits over the
algebraic closure of \(F\) after taking the component of the required
degree (Section~\ref{sec:notation}).

For the three-vertex path, Equation~\eqref{eq:graph} is
\(\sum_{i,j=1}^q x_iY_{ij}z_j\), the product of a row, a matrix and a
column. The formula bound is sharp in a regime that already suffices to
separate \(\VF\) from \(\VBP\). For a path on \(d\) vertices with \(q\)
states on each edge, bisection gives a formula of size \(q^{O(\log d)}\).
Theorem~\ref{thm:full-flat-graph} gives the matching
\(q^{\Omega(\log d)}\) bound for \(2\le d\le(\log q)^{1/3}\), and its
lower bound holds for every tree of that order. A short path already
suffices for the determinant bound in Theorem~\ref{thm:main}, because
its polynomial is a signed cofactor of a polynomial-size matrix of
affine forms.

For arbitrary graph families, these witnesses give a classification.
Consider a single polynomial bound on the size of \(Z_{G,q}\) in
\(|V(G)|+q\), uniformly over the family and all \(q\ge2\).
Formulas admit such a bound exactly when both the maximum degree and
the order of the largest connected component are bounded. For ABPs the
criterion is bounded cutwidth. For circuits it is bounded treewidth
and bounded maximum degree. Large graphs violating these conditions
either have a large table at one vertex or contain the corresponding
witnesses as minors (Theorem~\ref{thm:graph-classification}).

\needspace{7\baselineskip}
\subsection*{How the three lower bounds fit together}

The coefficients of \(Z_{G,q}\) record one constraint per edge. A
monomial selects one table entry at every vertex, and its coefficient
is one precisely when the entries agree on the state at both ends of
every edge. Thus the coefficient tensor consists of independent
equality constraints. The Euclidean price theorem of the first
paper~\cite{oaG}, based on Strassen's asymptotic spectrum~\cite{Strassen},
assigns a cost to sharing a label at chosen sites in Euclidean space.
We place the two ends of each edge at opposite unit vectors on an axis
of its own. The cost of
the graph tensor at these terminals is then the sum of the individual
edge costs.

A small computation expresses the coefficient tensor as a sum of terms
with simpler dependencies. Suppose a term has matrix rank \(R\) across
a cut separating the vertex tables into two groups. Its rank
factorization can be supplied by sharing a single \(R\)-valued label
between two auxiliary sites. For each term we may choose those sites
to suit its cut while delivering each edge-end index to its assigned
terminal. That terminal stays fixed across all summands. Moving one
supply site to shorten the total length of its output wires yields a
saving governed by the number of crossing edges. If \(R\)
is small and sufficiently many edges cross, the term costs less than
the graph tensor. A single Euclidean price bounds the sum, including
when the low-rank cut differs from term to term
(Section~\ref{sec:cuts}).

\begin{figure}[!hbp]
\centering
\begin{tikzpicture}[x=.48cm,y=.48cm,font=\small,
 ordinary/.style={draw=ink,line width=.7pt},
 crossing/.style={draw=rust,densely dashed,line width=1.15pt},
 chosen/.style={circle,draw=ink,fill=ink,inner sep=0pt,minimum size=4.3pt},
 outside/.style={circle,draw=ink,fill=white,inner sep=0pt,minimum size=4.3pt}]
\begin{scope}
  \foreach \i in {0,...,3} {
    \draw[ordinary] (\i,0)--(\i,3);
    \draw[ordinary] (0,\i)--(1,\i) (2,\i)--(3,\i);
    \draw[crossing] (1,\i)--(2,\i);
  }
  \foreach \i in {0,...,3} {
    \foreach \j in {0,1} \node[chosen] at (\j,\i) {};
    \foreach \j in {2,3} \node[outside] at (\j,\i) {};
  }
  \node at (1.5,-1.05) {Square grid};
\end{scope}
\begin{scope}[xshift=3.65cm]
  \coordinate (root) at (1.8,3);
  \coordinate (l) at (.6,1.5); \coordinate (r) at (3,1.5);
  \coordinate (ll) at (0,0); \coordinate (lr) at (1.2,0);
  \coordinate (rl) at (2.4,0); \coordinate (rr) at (3.6,0);
  \draw[crossing] (root)--(l);
  \draw[ordinary] (root)--(r) (l)--(ll) (l)--(lr) (r)--(rl) (r)--(rr);
  \foreach \v in {l,ll,lr} \node[chosen] at (\v) {};
  \foreach \v in {root,r,rl,rr} \node[outside] at (\v) {};
  \node at (1.8,-1.05) {Complete binary tree};
\end{scope}
\begin{scope}[xshift=7.35cm]
  \draw[ordinary] (0,1.5)--(2.25,1.5) (3,1.5)--(4.5,1.5);
  \draw[crossing] (2.25,1.5)--(3,1.5);
  \foreach \i in {0,...,3} \node[chosen] at (.75*\i,1.5) {};
  \foreach \i in {4,...,6} \node[outside] at (.75*\i,1.5) {};
  \node at (2.25,-1.05) {Path};
\end{scope}
\end{tikzpicture}
\caption{Three cuts between filled and open vertices. Four edges
cross the middle of this grid. On the binary tree a balanced split
can cross just one edge, so the ABP argument prescribes the size of
the smaller set more carefully. On the path, a proper nonempty prefix
of any size cuts just one edge. The formula argument combines splits
at several scales.}
\label{fig:witness-cuts}
\end{figure}


Figure~\ref{fig:witness-cuts} illustrates why the proofs proceed from
grids to trees to paths. Each model controls a different set of cuts.
\begin{itemize}[leftmargin=*,itemsep=.35ex]
\item \emph{Grids.} In a small circuit, decomposing at gates of
intermediate degree gives one low-rank cut per summand. Both sides
contain at least a third of the vertices. Every such cut of an
\(r\times r\) grid crosses \(\Omega(r)\) edges
(Section~\ref{sec:grid}).
\item \emph{Binary trees.} A tree can have a balanced cut consisting of
a single edge. An ABP gives more information. In each contributing term of the ABP path expansion, the variables
appear in some order. We can stop after any prescribed number \(k\) of
groups. There is a \(k\) for which
every \(k\)-vertex subset of the complete binary tree has many
boundary edges (Section~\ref{sec:bt}).
\item \emph{Paths.} For every nontrivial cardinality a path has a
prefix of that size with just one boundary edge. A formula, however,
splits the variable groups recursively. We choose the auxiliary
sites to charge its nested splits at several scales at once. The
mean-width formula bounds the cost of each label shared by the
leaves below a split (Section~\ref{sec:formula}).
\end{itemize}

The quantitative estimates first apply to small graphs with many states
per edge. In this regime the saving in each term exceeds the cost of
summing over terms. Projection transfers the bounds to the larger
witnesses in Theorem~\ref{thm:main}. The same transfer, now through
graph minors, gives the classification in
Section~\ref{sec:graph-classification}.

\subsection*{Relation to earlier work}

Valiant introduced $\VP$ and $\VNP$ to study efficient arithmetic
computation under polynomial projections~\cite{Valiant,Burgisser}.
The determinant is complete for $\VBP$ under such projections
\cite{Toda,MalodPortier,MahajanVinay}. Its formula complexity therefore
measures the cost of simulating branching programs by formulas.
The depth reduction of Valiant, Skyum, Berkowitz and Rackoff
\cite{VSBR} gives every family in $\VP$ formulas of
quasipolynomial size. A classical lower bound of Kalorkoti
\cite{Kalorkoti} gives $\Omega(n^3)$ for the order-$n$ determinant over
any field, counting variable and constant occurrences even in formulas
with division.

Many earlier separations use restrictions on the computation. Nisan's
coefficient matrices give exponential branching-program lower bounds
for explicit noncommutative polynomials~\cite{Nisan}. Raz proved
superpolynomial multilinear formula bounds for the determinant and
permanent, separated multilinear formulas from circuits, and related
formula size to tensor rank~\cite{RazMultilinear,RazFormula,RazRank}.
Dvir, Malod, Perifel and Yehudayoff separated multilinear formulas
from multilinear branching programs~\cite{DMPY}. For iterated matrix
multiplication itself, lower bounds are known for monotone formulas
\cite{Snir}, set-multilinear formulas~\cite{LST}, and constant-depth
circuits~\cite{LstDepth,Forbes,BDS}. The formulas and circuits of
Theorem~\ref{thm:main} may use arbitrary field constants and have
unbounded depth and nonmultilinear intermediate gates.

Graph structure also governs several related models. Hrube\v{s} and
Yehudayoff used the edge boundary of vertex sets of a prescribed size
in a binary tree to separate monotone branching programs from monotone
circuits~\cite{HrubesYehudayoff}. For graph homomorphism polynomials,
Komarath, Pandey and Rahul characterized monotone formula,
branching-program and circuit complexity by the treedepth, pathwidth
and treewidth of the pattern graph~\cite{KPR}. Homomorphism polynomials
also yield complete families for $\VP$~\cite{Durand} and separations
between symmetric models of arithmetic computation~\cite{dw,DPS}.
The classification in Theorem~\ref{thm:graph-classification} concerns
the state sum with an independent table of variables at each vertex.

For general circuits, an earlier bound for the order-$n$ permanent over
$\mathbb C$ was $\Omega(n^2\log\log n)$ with binary gates and individual
variable or scalar inputs~\cite[Chapter~5]{TenAdvances}. A companion paper~\cite{oaP} gives stronger circuit
bounds for the Hamiltonian cycle polynomial and, outside characteristic
two, the permanent. For the geometric study of determinantal
representations and their closures, see~\cite{MS}.
